% ============================================================
%  Beat by Bit — Documentação Técnica (Semestre 2)
%  Compilar com: pdflatex documentacao.tex
% ============================================================
\documentclass[12pt,a4paper]{article}

% ── Pacotes ────────────────────────────────────────────────
\usepackage[utf8]{inputenc}
\usepackage[T1]{fontenc}
\usepackage[brazilian]{babel}
\usepackage{geometry}
\usepackage{graphicx}
\usepackage{xcolor}
\usepackage{listings}
\usepackage{booktabs}
\usepackage{tabularx}
\usepackage{longtable}
\usepackage{array}
\usepackage{multirow}
\usepackage{amsmath}
\usepackage{hyperref}
\usepackage{tikz}
\usepackage{float}
\usepackage{fancyhdr}
\usepackage{titlesec}
\usepackage{enumitem}
\usepackage{mdframed}

\usetikzlibrary{shapes.geometric, arrows.meta, positioning, fit, backgrounds}

\geometry{
    top=2.5cm, bottom=2.5cm,
    left=2.5cm, right=2.5cm,
    headheight=14.5pt
}

% ── Cores do projeto ───────────────────────────────────────
\definecolor{fpgablue}{RGB}{0, 102, 204}
\definecolor{pygreen}{RGB}{0, 150, 70}
\definecolor{uartred}{RGB}{200, 50, 50}
\definecolor{codegray}{RGB}{245, 245, 245}
\definecolor{codeframe}{RGB}{180, 180, 180}
\definecolor{darkgray}{RGB}{60, 60, 60}
\definecolor{notecolor}{RGB}{255, 220, 0}

% ── Estilo de código ───────────────────────────────────────
\lstdefinestyle{verilog}{
    backgroundcolor=\color{codegray},
    frame=single,
    rulecolor=\color{codeframe},
    basicstyle=\ttfamily\footnotesize,
    keywordstyle=\color{fpgablue}\bfseries,
    commentstyle=\color{darkgray}\itshape,
    stringstyle=\color{uartred},
    numberstyle=\tiny\color{darkgray},
    numbers=left,
    numbersep=8pt,
    tabsize=4,
    breaklines=true,
    breakatwhitespace=true,
    showstringspaces=false,
    morekeywords={module, endmodule, input, output, wire, reg, always,
                  posedge, negedge, begin, end, if, else, case, endcase,
                  assign, parameter, localparam, integer, for, task,
                  endtask, function, endfunction},
    morecomment=[l]{//},
    morecomment=[s]{/*}{*/},
}

\lstdefinestyle{python}{
    backgroundcolor=\color{codegray},
    frame=single,
    rulecolor=\color{codeframe},
    basicstyle=\ttfamily\footnotesize,
    keywordstyle=\color{pygreen}\bfseries,
    commentstyle=\color{darkgray}\itshape,
    stringstyle=\color{uartred},
    numberstyle=\tiny\color{darkgray},
    numbers=left,
    numbersep=8pt,
    tabsize=4,
    breaklines=true,
    showstringspaces=false,
    language=Python,
}

\lstdefinestyle{shell}{
    backgroundcolor=\color{black!90},
    frame=none,
    basicstyle=\ttfamily\footnotesize\color{green!80!black},
    breaklines=true,
    showstringspaces=false,
}

% ── Cabeçalho ─────────────────────────────────────────────
\pagestyle{fancy}
\fancyhf{}
\fancyhead[L]{\small\textcolor{fpgablue}{\textbf{Beat by Bit}} — Documentação Técnica}
\fancyhead[R]{\small Semestre 2}
\fancyfoot[C]{\thepage}
\renewcommand{\headrulewidth}{0.4pt}

% ── Títulos ────────────────────────────────────────────────
\titleformat{\section}{\large\bfseries\color{fpgablue}}{}{0em}{\thesection\quad}
\titleformat{\subsection}{\normalsize\bfseries\color{darkgray}}{}{0em}{\thesubsection\quad}

\hypersetup{
    colorlinks=true,
    linkcolor=fpgablue,
    urlcolor=fpgablue,
}

% ── Documento ─────────────────────────────────────────────
\begin{document}

% ── Capa ──────────────────────────────────────────────────
\begin{titlepage}
    \centering
    \vspace*{2cm}

    \begin{tikzpicture}
        % Sombra escura (estilo 3D arcade)
        \node[font=\fontsize{45}{50}\sffamily\bfseries, text=black!80, xshift=3pt, yshift=-3pt] {BEAT BY BIT};
        % Texto principal em azul FPGA
        \node[font=\fontsize{45}{50}\sffamily\bfseries, text=fpgablue] {BEAT BY BIT};
        % Brilho (highlight neon) no topo/esquerda
        \node[font=\fontsize{45}{50}\sffamily\bfseries, text=cyan!30!white, xshift=-1pt, yshift=1pt] {BEAT BY BIT};
    \end{tikzpicture}\\[0.5cm]
    {\large Documentação Técnica — Semestre 2}\\[2cm]

    \begin{mdframed}[linecolor=fpgablue, linewidth=2pt, innertopmargin=12pt, innerbottommargin=12pt]
    \centering
    \textbf{Resumo das mudanças}\\[6pt]
    Migração da lógica do jogo do Python para a FPGA (Verilog).\\
    O Python passou a funcionar exclusivamente como \emph{renderer},\\
    recebendo o estado do jogo via UART em pacotes de 20 bytes a 60\,fps.
    \end{mdframed}

    \vspace{1.5cm}

    \begin{tabular}{ll}
        \textbf{Projeto anterior:} & \texttt{projeto\_sem1/} \\[4pt]
        \textbf{Projeto atual:}    & \texttt{projeto\_sem2/} \\[4pt]
        \textbf{FPGA:}             & Cyclone V (Intel/Altera) \\[4pt]
        \textbf{Clock:}            & 50\,MHz (assumido) \\[4pt]
        \textbf{UART:}             & 115200 baud, 8N1 \\[4pt]
        \textbf{Simulador:}        & ModelSim / Icarus Verilog \\[4pt]
        \textbf{Renderer:}         & Python 3 + Pygame \\
    \end{tabular}

    \vfill
    {\small Laboratório de Circuitos Digitais}
\end{titlepage}

\tableofcontents
\newpage

% ═══════════════════════════════════════════════════════════
\section{Visão Geral}
% ═══════════════════════════════════════════════════════════

\subsection{Objetivo}

O projeto \textbf{Beat by Bit} é um jogo de ritmo implementado em FPGA. No semestre 1, toda a lógica do jogo — geração de notas, detecção de acerto, pontuação, temporizadores e condições de vitória — estava implementada em Python. A FPGA apenas transmitia o estado dos botões físicos para o computador.

No semestre 2, a lógica foi \textbf{completamente migrada para o Verilog}, respeitando o requisito do trabalho: \emph{a FPGA deve ser o cérebro do jogo}. O Python passou a ser um \emph{renderer} puro — só recebe pacotes via UART e desenha a tela.

\subsection{Divisão de responsabilidades}

\begin{table}[H]
\centering
\begin{tabularx}{\textwidth}{Xcc}
\toprule
\textbf{Funcionalidade} & \textbf{Sem.~1} & \textbf{Sem.~2} \\
\midrule
Geração de notas (spawn)         & Python & \textbf{FPGA} \\
Posição e movimento das notas    & Python & \textbf{FPGA} \\
Detecção de acerto (hit window)  & Python & \textbf{FPGA} \\
Pontuação (score / combo)        & Python & \textbf{FPGA} \\
Contador de erros (misses)       & Python & \textbf{FPGA} \\
Timer de jogo (60\,s)            & Python & \textbf{FPGA} \\
Countdown (3\,s)                 & Parcial & \textbf{FPGA} \\
Condição de vitória/derrota      & Python & \textbf{FPGA} \\
Máquina de estados (FSM)         & \textbf{FPGA} & \textbf{FPGA} \\
Debounce e detecção de borda     & \textbf{FPGA} & \textbf{FPGA} \\
Renderização gráfica             & Python & Python \\
\bottomrule
\end{tabularx}
\caption{Comparação de responsabilidades entre os semestres.}
\end{table}

% ═══════════════════════════════════════════════════════════
\section{Arquitetura do Sistema}
% ═══════════════════════════════════════════════════════════

\subsection{Diagrama de blocos}

\begin{figure}[H]
\centering
\begin{tikzpicture}[
    node distance=1.2cm and 1.8cm,
    box/.style={draw, rounded corners, minimum width=2.8cm, minimum height=0.9cm,
                fill=codegray, font=\small},
    newbox/.style={draw, rounded corners, minimum width=2.8cm, minimum height=0.9cm,
                   fill=fpgablue!15, draw=fpgablue, font=\small},
    label/.style={font=\footnotesize\itshape},
    arr/.style={-{Latex}, thick},
]

% Bloco FPGA (grande)
\node[box] (uc)  at (0,0)          {\texttt{unidade\_controle}};
\node[box] (fd)  at (0,-1.8)       {\texttt{fluxo\_dados}};
\node[newbox] (nt0) at (-3.5,-3.2) {\texttt{note\_track[0]}};
\node[newbox] (nt1) at (-1.0,-3.2) {\texttt{note\_track[1]}};
\node[newbox] (nt2) at ( 1.5,-3.2) {\texttt{note\_track[2]}};
\node[newbox] (nt3) at ( 4.0,-3.2) {\texttt{note\_track[3]}};
\node[newbox] (lfsr) at (-3.5,-4.8){\texttt{lfsr8}};
\node[newbox] (pkt) at (4.5,-0.9)  {\texttt{packet\_sender}};
\node[newbox] (utx) at (4.5,-2.4)  {\texttt{uart\_tx}};
\node[box] (hex) at (0, 1.5)       {\texttt{hexa7seg ×2}};

% PC
\node[draw, fill=pygreen!15, draw=pygreen, rounded corners,
      minimum width=2.8cm, minimum height=0.9cm, font=\small]
      (py) at (4.5,-4.2) {\textbf{Python} (renderer)};

% Entradas
\node[label, left=1.2cm of uc] (btn) {botoes[4:0]};
\node[label, left=1.2cm of fd] (clk) {clock, reset};

% Setas
\draw[arr] (btn.east) -- (uc.west);
\draw[arr] (clk.east) -- (fd.west);
\draw[arr] (uc) -- (fd) node[midway, right, label]{controle};
\draw[arr] (fd.south) -- ++(0,-0.3) -| (nt0.north);
\draw[arr] (fd.south) -- ++(0,-0.3) -| (nt1.north);
\draw[arr] (fd.south) -- ++(0,-0.3) -| (nt2.north);
\draw[arr] (fd.south) -- ++(0,-0.3) -| (nt3.north);
\draw[arr] (lfsr.north) -- (nt0.south);
\draw[arr, color=fpgablue] (fd.east) -- (pkt.west)
      node[midway, above, label, color=fpgablue]{estado};
\draw[arr, color=fpgablue] (pkt) -- (utx)
      node[midway, right, label, color=fpgablue]{bytes};
\draw[arr, color=uartred] (utx) -- (py)
      node[midway, right, label, color=uartred]{UART};
\draw[arr] (uc.north) -- (hex.south);

% Bounding box FPGA
\begin{scope}[on background layer]
\node[draw=fpgablue, dashed, rounded corners=8pt, fill=fpgablue!5,
      fit=(uc)(fd)(nt0)(nt3)(lfsr)(hex)(pkt)(utx),
      label=above:{\textbf{\textcolor{fpgablue}{FPGA — Cyclone V}}}] {};
\end{scope}

\end{tikzpicture}
\caption{Diagrama de blocos do sistema. Módulos novos destacados em azul.}
\end{figure}

\subsection{Hierarquia de módulos}

\begin{table}[H]
\centering
\begin{tabularx}{\textwidth}{lXl}
\toprule
\textbf{Arquivo} & \textbf{Descrição} & \textbf{Status} \\
\midrule
\texttt{projeto.v}          & Top-level; conecta todos os módulos  & Atualizado \\
\texttt{unidade\_controle.v}& FSM principal do jogo (6 estados)    & Atualizado \\
\texttt{fluxo\_dados.v}     & Lógica completa do jogo              & \textcolor{fpgablue}{\textbf{Reescrito}} \\
\texttt{uart\_tx.v}         & Transmissor UART 8N1                 & \textcolor{fpgablue}{\textbf{Novo}} \\
\texttt{lfsr8.v}            & LFSR 8 bits para spawn aleatório     & \textcolor{fpgablue}{\textbf{Novo}} \\
\texttt{note\_track.v}      & Gerenciador de notas por trilha      & \textcolor{fpgablue}{\textbf{Novo}} \\
\texttt{packet\_sender.v}   & Montagem e envio do pacote UART      & \textcolor{fpgablue}{\textbf{Novo}} \\
\texttt{contador\_m.v}      & Contador modulo-M parametrizável     & Reutilizado \\
\texttt{edge\_detector.v}   & Detector de borda de subida          & Reutilizado \\
\texttt{registrador\_4.v}   & Registrador 4 bits com enable        & Reutilizado \\
\texttt{hexa7seg.v}         & Decodificador para display 7 seg.    & Reutilizado \\
\texttt{interface.py}       & Renderer Python (Pygame)             & \textcolor{pygreen}{\textbf{Reescrito}} \\
\bottomrule
\end{tabularx}
\caption{Hierarquia e status dos arquivos do projeto.}
\end{table}

% ═══════════════════════════════════════════════════════════
\section{Protocolo UART}
% ═══════════════════════════════════════════════════════════

\subsection{Configuração}

\begin{itemize}
    \item \textbf{Baud rate:} 115\,200 bps
    \item \textbf{Formato:} 8N1 (8 bits de dados, sem paridade, 1 stop bit)
    \item \textbf{Direção:} FPGA $\rightarrow$ PC (unidirecional)
    \item \textbf{Taxa de envio:} 1 pacote por frame ($\approx$\,60\,fps)
    \item \textbf{Tempo por pacote:} $20 \times \frac{10}{115200} \approx 1{,}74\,\text{ms}$
\end{itemize}

\subsection{Formato do pacote (20 bytes)}

\begin{table}[H]
\centering
\begin{tabular}{cclc}
\toprule
\textbf{Byte(s)} & \textbf{Nome} & \textbf{Descrição} & \textbf{Valores} \\
\midrule
0  & \texttt{HEADER}     & Marcador de início                   & \texttt{0x55} \\
1  & \texttt{state}      & Estado da FSM                        & 0--5 \\
2  & \texttt{countdown}  & Segundos restantes no countdown      & 3, 2, 1, 0 \\
3  & \texttt{t0n0}       & Trilha 0 — age da nota 0             & 0--150, \texttt{0xFF} \\
4  & \texttt{t0n1}       & Trilha 0 — age da nota 1             & 0--150, \texttt{0xFF} \\
5  & \texttt{t0n2}       & Trilha 0 — age da nota 2             & 0--150, \texttt{0xFF} \\
6--8  & \texttt{t1n0--t1n2} & Trilha 1 (idem)                   & 0--150, \texttt{0xFF} \\
9--11 & \texttt{t2n0--t2n2} & Trilha 2 (idem)                   & 0--150, \texttt{0xFF} \\
12--14& \texttt{t3n0--t3n2} & Trilha 3 (idem)                   & 0--150, \texttt{0xFF} \\
15 & \texttt{score\_h}   & Pontuação (byte alto)                & 0--255 \\
16 & \texttt{score\_l}   & Pontuação (byte baixo)               & 0--255 \\
17 & \texttt{misses}     & Contador de erros                    & 0--10 \\
18 & \texttt{combo}      & Combo atual                          & 0--255 \\
19 & \texttt{FOOTER}     & Marcador de fim                      & \texttt{0xAA} \\
\bottomrule
\end{tabular}
\caption{Estrutura do pacote UART de 20 bytes. \texttt{0xFF} indica slot vazio.}
\end{table}

\subsection{Codificação de posição das notas}

Cada nota é armazenada na FPGA como um contador de \emph{age} (quadros desde o spawn). O Python converte em posição $Y$ na tela pela fórmula:

\begin{equation}
    Y = \mathit{age} \times \mathit{NOTE\_SPEED} \qquad (\mathit{NOTE\_SPEED} = 4\ \text{px/frame})
\end{equation}

\begin{table}[H]
\centering
\begin{tabular}{ccl}
\toprule
\textbf{Age} & \textbf{Y (px)} & \textbf{Significado} \\
\midrule
0   & 0   & Nota acabou de aparecer (topo da tela) \\
119 & 476 & Borda inferior da hit zone \\
130 & 520 & Centro da hit zone \\
141 & 564 & Borda superior da hit zone \\
151 & 604 & Nota saiu da tela $\rightarrow$ miss \\
\texttt{0xFF} & --- & Slot vazio (sem nota) \\
\bottomrule
\end{tabular}
\caption{Mapeamento age $\rightarrow$ posição Y.}
\end{table}

\subsection{Estados da FSM (byte 1 do pacote)}

\begin{center}
\begin{tabular}{cll}
\toprule
\textbf{Valor} & \textbf{Estado} & \textbf{Descrição} \\
\midrule
0 & \texttt{IDLE}      & Aguardando pressionar START \\
1 & \texttt{COUNTDOWN} & Contagem regressiva 3-2-1 \\
2 & \texttt{PLAY}      & Jogo em andamento \\
3 & \texttt{PAUSE}     & Jogo pausado \\
4 & \texttt{END\_WIN}  & Vitória (60\,s concluídos) \\
5 & \texttt{END\_LOSE} & Derrota (10 erros) \\
\bottomrule
\end{tabular}
\end{center}

% ═══════════════════════════════════════════════════════════
\section{Módulos Verilog Novos}
% ═══════════════════════════════════════════════════════════

\subsection{\texttt{uart\_tx.v} — Transmissor UART}

\textbf{Função:} Serializa um byte de 8 bits no formato 8N1 (1 start bit + 8 data bits + 1 stop bit).

\textbf{Parâmetro:}
\begin{itemize}
    \item \texttt{CLKS\_PER\_BIT} $= \lfloor f_{clk} / \mathit{baud} \rfloor = \lfloor 50\,000\,000 / 115\,200 \rfloor = 434$
\end{itemize}

\textbf{Portas:}
\begin{center}
\begin{tabular}{lll}
\toprule
\textbf{Sinal} & \textbf{Direção} & \textbf{Descrição} \\
\midrule
\texttt{clock}     & entrada & Clock do sistema \\
\texttt{reset}     & entrada & Reset assíncrono ativo alto \\
\texttt{tx\_data[7:0]} & entrada & Byte a transmitir \\
\texttt{tx\_valid} & entrada & Pulso de 1 ciclo para iniciar transmissão \\
\texttt{tx\_serial}& saída   & Sinal serial de saída (pino UART TX) \\
\texttt{tx\_ready} & saída   & \texttt{1} quando pronto para novo byte \\
\bottomrule
\end{tabular}
\end{center}

\textbf{Diagrama de estados:}
\begin{center}
\begin{tikzpicture}[
    state/.style={draw, circle, minimum size=1.8cm, font=\small},
    arr/.style={-{Latex}, thick},
]
\node[state, fill=fpgablue!15] (idle)  at (0,0)   {\texttt{IDLE}};
\node[state] (start) at (3,0)  {\texttt{START}};
\node[state] (data)  at (6,0)  {\texttt{DATA}};
\node[state] (stop)  at (9,0)  {\texttt{STOP}};
\draw[arr] (idle)  -- (start) node[midway, above, font=\footnotesize]{\texttt{tx\_valid}};
\draw[arr] (start) -- (data)  node[midway, above, font=\footnotesize]{fim bit};
\draw[arr] (data)  -- (stop)  node[midway, above, font=\footnotesize]{8 bits};
\draw[arr] (stop)  to[bend left=40] (idle)  node[midway, below, font=\footnotesize]{fim bit};
\end{tikzpicture}
\end{center}

\subsection{\texttt{lfsr8.v} — Gerador Pseudoaleatório}

\textbf{Função:} LFSR (\emph{Linear Feedback Shift Register}) de 8 bits para gerar spawn pseudoaleatório das notas.

\textbf{Polinômio:} $x^8 + x^6 + x^5 + x^4 + 1$ \quad (taps nos bits 8, 6, 5, 4)

\textbf{Propriedades:}
\begin{itemize}
    \item Período máximo de 255 estados (nunca retorna ao estado zero)
    \item Avança 1 passo a cada pulso de \texttt{enable} (\texttt{frame\_tick})
    \item Semente inicial: \texttt{0xAC}
\end{itemize}

\textbf{Uso no spawn:} O \texttt{fluxo\_dados} usa 2 bits do LFSR por trilha. Uma nota é spawnada se os 2 bits forem \texttt{00} (probabilidade $\approx 25\%$ por trilha por frame).

\begin{lstlisting}[style=verilog, caption={Lógica de feedback do LFSR8.}]
wire feedback = lfsr_reg[7] ^ lfsr_reg[5] ^ lfsr_reg[4] ^ lfsr_reg[3];
lfsr_reg <= {lfsr_reg[6:0], feedback};  // shift left com feedback
\end{lstlisting}

\subsection{\texttt{note\_track.v} — Gerenciador de Notas por Trilha}

\textbf{Função:} Mantém até 3 notas simultâneas numa trilha. Cada nota é representada por um contador de age (quadros desde o spawn).

\textbf{Constantes de jogo:}
\begin{center}
\begin{tabular}{lll}
\toprule
\textbf{Constante} & \textbf{Valor} & \textbf{Significado} \\
\midrule
\texttt{HIT\_CENTER} & 130 & Age quando a nota está no centro da hit zone ($Y=520$) \\
\texttt{HIT\_WINDOW} & 11  & Janela de tolerância $\pm11$ frames ($\pm44$\,px) \\
\texttt{MISS\_AGE}   & 151 & Age em que a nota sai da tela sem ser apertada \\
\texttt{EMPTY}       & \texttt{0xFF} & Slot vazio \\
\bottomrule
\end{tabular}
\end{center}

\textbf{Lógica por \texttt{frame\_tick}:}
\begin{enumerate}
    \item Para todos os slots válidos: \texttt{age++}
    \item Se \texttt{age >= 151}: marca \texttt{escape\_pulse=1}, libera slot
    \item Se \texttt{spawn\_en} e há slot livre: cria nota com \texttt{age=0}
\end{enumerate}

\textbf{Lógica por \texttt{btn\_press}:}
\begin{enumerate}
    \item Procura nota com age dentro da janela $[119,\,141]$
    \item Se encontrada: \texttt{hit\_pulse=1}, remove a nota
    \item Caso contrário: \texttt{bad\_press=1}
\end{enumerate}

\textbf{Saídas de pulso} (duram exatamente 1 ciclo de clock):
\begin{center}
\begin{tabular}{ll}
\toprule
\textbf{Sinal} & \textbf{Significado} \\
\midrule
\texttt{hit\_pulse}    & Botão pressionado dentro da janela → acerto \\
\texttt{escape\_pulse} & Nota passou sem ser apertada → miss \\
\texttt{bad\_press}    & Botão pressionado fora da janela → penalidade \\
\bottomrule
\end{tabular}
\end{center}

\subsection{\texttt{packet\_sender.v} — Montagem e Envio do Pacote}

\textbf{Função:} A cada \texttt{frame\_tick}, captura o estado completo do jogo, monta o pacote de 20 bytes e envia byte a byte via \texttt{uart\_tx}.

\textbf{Máquina de estados interna:}

\begin{center}
\begin{tikzpicture}[
    state/.style={draw, rounded corners, minimum width=2cm, minimum height=0.8cm,
                  fill=codegray, font=\small},
    arr/.style={-{Latex}, thick},
    node distance=0.8cm and 2cm,
]
\node[state, fill=fpgablue!15] (wait) {WAIT};
\node[state, right=of wait] (load) {LOAD};
\node[state, right=of load] (send) {SEND};
\node[state, right=of send] (done) {DONE};

\draw[arr] (wait) -- (load) node[midway, above, font=\tiny]{\texttt{frame\_tick}};
\draw[arr] (load) -- (send) node[midway, above, font=\tiny]{\texttt{tx\_ready}};
\draw[arr] (send) to[bend right=20] (load)
    node[midway, below, font=\tiny]{próx. byte};
\draw[arr] (send) -- (done) node[midway, above, font=\tiny]{byte 19};
\draw[arr] (done) to[bend left=40] (wait)
    node[midway, below, font=\tiny]{\texttt{tx\_ready}};
\end{tikzpicture}
\end{center}

O estado \texttt{WAIT} captura (\emph{latch}) todos os sinais de entrada no momento do \texttt{frame\_tick}, garantindo consistência do pacote. O envio acontece nos estados seguintes sem risco de leitura de dados em transição.

\subsection{\texttt{fluxo\_dados.v} — Lógica Completa do Jogo}

Este módulo centraliza toda a lógica do jogo. Instancia os submódulos e mantém os registradores de estado.

\textbf{Temporizadores:}

\begin{table}[H]
\centering
\begin{tabular}{llll}
\toprule
\textbf{Temporizador} & \textbf{M} & \textbf{N (bits)} & \textbf{Duração} \\
\midrule
Frame tick (\texttt{u\_frame\_counter}) & 833\,333 & 20 & $\approx 1/60\,\text{s}$ \\
Countdown (\texttt{u\_countdown})       & 180      & 8  & 3 segundos (180 frames) \\
Timer de jogo (\texttt{u\_game\_timer}) & 3\,600   & 12 & 60 segundos \\
\bottomrule
\end{tabular}
\caption{Temporizadores baseados em \texttt{contador\_m} instanciados no \texttt{fluxo\_dados}. Clock de 50\,MHz assumido.}
\end{table}

\textbf{Debounce dos botões:} Cada entrada física em \texttt{botoes\_raw[4:0]} agora passa por um módulo \texttt{debounce\_pulse}. O primeiro toque gera um pulso de 1 ciclo logo na primeira borda de subida sincronizada, e novas transições ficam bloqueadas por 180\,ms. Para clock de 50\,MHz, isso corresponde a 9\,000\,000 ciclos. A política foi aplicada tanto às 4 trilhas quanto ao botão \texttt{START/PAUSE}; como consequência, toques sucessivos na mesma trilha com menos de 180\,ms entre si são descartados.

\textbf{Contadores de placar:}

\begin{lstlisting}[style=verilog, caption={Lógica dos contadores de pontuação no \texttt{fluxo\_dados.v}.}]
// Score (+100 acerto, -50 bad_press, minimo 0)
always @(posedge clock or posedge reset)
    if (reset || limpaR)    score_reg <= 0;
    else if (any_hit)       score_reg <= score_reg + 16'd100;
    else if (any_bad)       score_reg <= (score_reg >= 16'd50) ?
                                         score_reg - 16'd50 : 16'd0;

// Misses (nota escapa -> +1; perdeu quando >= 10)
always @(posedge clock or posedge reset)
    if (reset || limpaR)              miss_reg <= 0;
    else if (any_escape && game_active) miss_reg <= miss_reg + 1;

assign perdeu = (miss_reg >= 8'd10);

// Combo (acerto -> +1; erro -> reset)
always @(posedge clock or posedge reset)
    if (reset || limpaR)             combo_reg <= 0;
    else if (any_hit)                combo_reg <= combo_reg + 1;
    else if (any_bad || any_escape)  combo_reg <= 0;
\end{lstlisting}

\subsection{\texttt{unidade\_controle.v} — FSM Principal}

A FSM permaneceu com os mesmos 6 estados. A alteração foi a adição da saída \texttt{game\_active} (HIGH somente no estado \texttt{PLAY}), que habilita o movimento das notas e os contadores de placar no \texttt{fluxo\_dados}.

\begin{center}
\begin{tikzpicture}[
    state/.style={draw, circle, minimum size=1.6cm, font=\footnotesize, align=center},
    arr/.style={-{Latex}},
    node distance=1.5cm,
]
\node[state, fill=codegray]     (idle) {IDLE\\0};
\node[state, right=2cm of idle] (cd)   {COUNT\\DOWN\\1};
\node[state, right=2cm of cd]   (play) {PLAY\\2};
\node[state, right=2cm of play] (pause){PAUSE\\3};
\node[state, below=1.8cm of play] (win){END\\WIN\\4};
\node[state, below=1.8cm of cd]   (lose){END\\LOSE\\5};

\draw[arr] (idle) -- (cd)   node[midway, above, font=\tiny]{START};
\draw[arr] (cd)   -- (play) node[midway, above, font=\tiny]{fim\_contagem};
\draw[arr] (play) to[bend right=20] (pause) node[midway, above, font=\tiny]{START};
\draw[arr] (pause) to[bend right=20] (play) node[midway, below, font=\tiny]{START};
\draw[arr] (play) -- (win)  node[midway, right, font=\tiny]{fim\_tempo};
\draw[arr] (play) -- (lose) node[midway, left, font=\tiny]{perdeu};
\draw[arr] (win)  to[bend left=30] (idle) node[near start, right, font=\tiny]{START};
\draw[arr] (lose) to[bend right=30] (idle) node[near start, left, font=\tiny]{START};
\end{tikzpicture}
\end{center}

% ═══════════════════════════════════════════════════════════
\section{Integração Verilog-Python (Renderer)}
% ═══════════════════════════════════════════════════════════

O script \texttt{interface.py} atua como o motor audiovisual (renderer) do sistema. O paradigma arquitetural adotado define que \textbf{toda a lógica de estado e regras de negócio residem na FPGA (Verilog)}, enquanto o \textbf{Python atua passivamente}, refletindo o estado do hardware na tela.

\subsection{Dinâmica de Comunicação (Produtor-Consumidor)}

A integração segue um modelo de produtor-consumidor em tempo real:
\begin{itemize}
    \item \textbf{Produtor (FPGA):} A cada borda de subida do \texttt{frame\_tick} ($\approx$ 60 vezes por segundo), o módulo \texttt{packet\_sender.v} faz um \emph{latch} (captura) de todas as variáveis críticas do jogo (score, combo, age das notas, estado da FSM). Esses dados são serializados e enviados via \texttt{uart\_tx} em um pacote binário de 20 bytes.
    \item \textbf{Consumidor (Python):} O PC recebe os sinais através da porta USB (usando o padrão UART). O script Python utiliza a biblioteca \texttt{pyserial} dentro de uma \emph{thread} dedicada (\texttt{SerialReader}) para drenar o buffer do sistema operacional continuamente, evitando gargalos ou travamentos no loop de renderização gráfica do Pygame.
\end{itemize}

\subsection{Sincronização e Reconstrução do Estado}

Como a comunicação UART assíncrona está sujeita a atrasos ou perda de bits caso o cabo sofra interferência, a integração garante a integridade dos dados através de \textbf{marcadores de fronteira}:
\begin{enumerate}
    \item A função \texttt{find\_packet} varre o buffer serial em busca do byte de cabeçalho (\texttt{0x55}).
    \item Ao encontrar, verifica se o byte na posição $+19$ é o terminador (\texttt{0xAA}).
    \item Se a condição for satisfeita, o pacote é validado e a função \texttt{parse\_packet} extrai os 18 bytes intermediários para reconstruir a matriz de notas e as variáveis numéricas de HUD.
    \item Caso o pacote esteja fragmentado, o Python descarta os bytes iniciais e busca o próximo alinhamento válido, garantindo ressincronização automática quase instantânea.
\end{enumerate}

\begin{lstlisting}[style=python, caption={Trecho do Parser verificando a integridade do pacote UART.}]
def parse_packet(data, offset):
    if data[offset] != 0x55 or data[offset + 19] != 0xAA:
        return None  # Pacote invalido ou incompleto
    state     = data[offset + 1] & 0x0F
    # ... (extracao de notas e placar)
    return {'state': state, 'notes': notes, 'score': score, ...}
\end{lstlisting}

\subsection{Mapeamento Visual e Sonoro Indireto}

Sendo um \emph{renderer} passivo, o Python não calcula a física do jogo nem verifica se um botão foi apertado no momento certo. Ele infere os efeitos visuais de forma inteligente comparando as diferenças entre os pacotes recebidos:
\begin{itemize}
    \item \textbf{Animação de Notas:} A coordenada $Y$ exata das notas na tela não é calculada por fórmulas de movimento no Python, mas sim derivada diretamente da "idade" da nota (\texttt{age}) enviada pelo Verilog multiplicada por uma constante.
    \item \textbf{Feedback Visual (Hits e Misses):} O Python guarda o estado das notas do frame anterior. Se ele percebe que uma nota desapareceu da lista (comparando os conjuntos \texttt{prev\_notes} e \texttt{curr\_notes}) e, simultaneamente, detecta que o contador de \texttt{score} ou \texttt{combo} enviado pela FPGA aumentou, ele sabe que o usuário acertou e desenha na tela a animação de acerto (partículas explodindo e luzes neon nos botões arcade). Se a nota sumiu abaixo da tela e os erros (\texttt{misses}) subiram, ele pisca um flash vermelho de erro nas bordas da trilha correspondente.
    \item \textbf{Áudio e Interface:} O byte numérico de estado (\texttt{STATE\_IDLE}, \texttt{STATE\_PLAY}, \texttt{STATE\_SONG\_SELECT}, etc.) enviado pela máquina de estados da FPGA controla qual tela é exibida e dita os comandos do reprodutor de música (iniciando o áudio no Play, aplicando o \texttt{unpause} correto ao voltar de um Pause, ou parando completamente).
\end{itemize}

% ═══════════════════════════════════════════════════════════
\section{Como Fazer Funcionar na Prática}
% ═══════════════════════════════════════════════════════════

\subsection{Requisitos}

\begin{itemize}
    \item \textbf{Hardware:} Placa FPGA com Cyclone V (ex: DE1-SoC, DE2-115)
    \item \textbf{Software FPGA:} Intel Quartus Prime (versão Lite é suficiente)
    \item \textbf{Software PC:} Python 3.8+, Pygame, pyserial
    \item \textbf{Cabo:} USB-A para mini/micro-USB (conexão UART da placa)
\end{itemize}

\subsection{Passo 1: Configurar o Projeto no Quartus}

\begin{enumerate}
    \item Abra o Quartus Prime e crie um novo projeto ou abra o \texttt{beat\_by\_bit.qar} existente.
    \item \textbf{Adicione os novos arquivos Verilog} ao projeto (\textit{Project} $\rightarrow$ \textit{Add/Remove Files}):
    \begin{itemize}
        \item \texttt{uart\_tx.v}
        \item \texttt{lfsr8.v}
        \item \texttt{note\_track.v}
        \item \texttt{packet\_sender.v}
    \end{itemize}
    \item Certifique-se de que \texttt{projeto.v} está definido como \textbf{Top-Level Entity}.
    \item Verifique se o dispositivo está configurado para o Cyclone V correto (\textit{Assignments} $\rightarrow$ \textit{Device}).
\end{enumerate}

\subsection{Passo 2: Ajustar o Clock}

O parâmetro \texttt{M} do contador de frame é calculado por:

\begin{equation}
    M_{\mathit{frame}} = \left\lfloor \frac{f_{clk}}{60} \right\rfloor
\end{equation}

\begin{center}
\begin{tabular}{ccc}
\toprule
\textbf{Clock da placa} & \textbf{M (frame tick)} & \textbf{N (bits)} \\
\midrule
50\,MHz  & 833\,333 & 20 \\
100\,MHz & 1\,666\,667 & 21 \\
25\,MHz  & 416\,667  & 19 \\
\bottomrule
\end{tabular}
\end{center}

Para verificar a frequência do clock da sua placa, consulte o manual ou o arquivo \texttt{c5\_pin\_model\_dump.txt}. Edite \texttt{fluxo\_dados.v} linha correspondente ao \texttt{u\_frame\_counter}:

\begin{lstlisting}[style=verilog]
contador_m #(.M(833_333), .N(20)) u_frame_counter  // 50 MHz
\end{lstlisting}

Da mesma forma, o parâmetro do \texttt{uart\_tx} deve ser ajustado:

\begin{equation}
    \mathtt{CLKS\_PER\_BIT} = \left\lfloor \frac{f_{clk}}{115200} \right\rfloor
\end{equation}

Edite \texttt{packet\_sender.v}:
\begin{lstlisting}[style=verilog]
uart_tx #(.CLKS_PER_BIT(434)) u_tx  // 50 MHz / 115200
\end{lstlisting}

\subsection{Passo 3: Atribuir os Pinos (Pin Planner)}

Abra o \textit{Pin Planner} (\textit{Assignments} $\rightarrow$ \textit{Pin Planner}) e atribua:

\begin{table}[H]
\centering
\begin{tabular}{llll}
\toprule
\textbf{Sinal no Verilog} & \textbf{Pino típico DE1-SoC} & \textbf{Descrição} \\
\midrule
\texttt{clock}         & PIN\_AF14      & Clock de 50\,MHz \\
\texttt{reset}         & KEY[0]         & Botão de reset \\
\texttt{botoes[0]}     & KEY[1]         & Trilha amarela (D) \\
\texttt{botoes[1]}     & Botão físico   & Trilha azul (F) \\
\texttt{botoes[2]}     & Botão físico   & Trilha verde (J) \\
\texttt{botoes[3]}     & Botão físico   & Trilha vermelha (K) \\
\texttt{botoes[4]}     & KEY[3]         & START / PAUSE \\
\texttt{uart\_tx\_out} & PIN do UART TX & TX da UART USB \\
\texttt{db\_estado}    & HEX0[6:0]      & Display 7 segmentos \\
\texttt{db\_jogada}    & HEX1[6:0]      & Display 7 segmentos \\
\texttt{leds\_pulsos}  & LEDR[4:0]      & LEDs de debug \\
\bottomrule
\end{tabular}
\caption{Atribuição de pinos para placa DE1-SoC. Ajuste conforme o manual da sua placa.}
\end{table}

\begin{mdframed}[linecolor=uartred, linewidth=1.5pt, innertopmargin=8pt, innerbottommargin=8pt]
\textbf{Atenção — Pino UART TX:} O pino de TX da UART USB varia conforme a placa. Na DE1-SoC é geralmente ligado ao \textit{HPS UART}. Consulte o \textit{User Manual} da placa e o arquivo \texttt{c5\_pin\_model\_dump.txt} para localizar o pino correto. É possível também usar qualquer GPIO livre e conectar um adaptador USB-Serial externo.
\end{mdframed}

\subsection{Passo 4: Compilar e Gravar}

\begin{lstlisting}[style=shell]
# No Quartus: Processing -> Start Compilation
# Verificar: 0 erros criticos no Flow Summary

# Gravar na FPGA: Tools -> Programmer
# Selecionar arquivo .sof gerado
# Clicar em "Start"
\end{lstlisting}

\subsection{Passo 5: Instalar Dependências Python}

\begin{lstlisting}[style=shell]
cd /home/zocal/labdig_files/projeto_sem2

# Ativar o ambiente virtual (se existir)
source env/bin/activate

# Instalar dependencias (caso necessario)
pip install pygame pyserial
\end{lstlisting}

\subsection{Passo 6: Executar o Renderer}

\begin{lstlisting}[style=shell]
# Verificar a porta serial da FPGA
ls /dev/ttyUSB*     # Linux
# ou
ls /dev/cu.usbserial*  # macOS

# Executar
python interface.py --port /dev/ttyUSB0 --baud 115200
\end{lstlisting}

\textbf{Parâmetros de linha de comando:}
\begin{center}
\begin{tabular}{lll}
\toprule
\textbf{Argumento} & \textbf{Padrão} & \textbf{Descrição} \\
\midrule
\texttt{--port} & \texttt{/dev/ttyUSB0} & Porta serial da FPGA \\
\texttt{--baud} & \texttt{115200}       & Baud rate \\
\bottomrule
\end{tabular}
\end{center}

\subsection{Indicadores visuais de status}

No canto superior direito da tela, o renderer exibe:
\begin{itemize}
    \item \textcolor{pygreen}{\textbf{FPGA ON}} — conexão UART estabelecida
    \item \textcolor{darkgray}{\textbf{FPGA OFF}} — sem conexão (modo demo, tela fica na última posição recebida)
\end{itemize}

\subsection{Solução de Problemas}

\begin{table}[H]
\centering
\begin{tabularx}{\textwidth}{lX}
\toprule
\textbf{Sintoma} & \textbf{Solução} \\
\midrule
Tela permanece em IDLE &
  Verificar se a FPGA está gravada e transmitindo. Testar com \texttt{minicom} ou \texttt{screen /dev/ttyUSB0 115200}. \\[6pt]
Notas aparecem em posições erradas &
  Verificar se o clock da FPGA é 50\,MHz. Se diferente, ajustar \texttt{M} no \texttt{fluxo\_dados.v}. \\[6pt]
Jogo muito lento ou rápido &
  O \texttt{NOTE\_SPEED} no Python deve ser igual ao usado no cálculo Verilog (ambos = 4). \\[6pt]
\texttt{/dev/ttyUSB0} não existe &
  Verificar driver USB-UART (\texttt{lsusb}), adicionar usuário ao grupo serial: \texttt{sudo usermod -aG dialout \$USER}. \\[6pt]
Permissão negada na porta serial &
  Executar \texttt{sudo chmod 666 /dev/ttyUSB0} ou adicionar ao grupo \texttt{dialout}. \\[6pt]
Pacotes corrompidos / estado errado &
  O \texttt{find\_packet} se ressincroniza automaticamente. Se persistir, verificar conexão do cabo e gravar novamente na FPGA. \\
\bottomrule
\end{tabularx}
\caption{Solução de problemas comuns.}
\end{table}

% ═══════════════════════════════════════════════════════════
\section{Simulação e Testes}
% ═══════════════════════════════════════════════════════════

\subsection{Testbenches disponíveis}

\begin{table}[H]
\centering
\begin{tabularx}{\textwidth}{llX}
\toprule
\textbf{Arquivo} & \textbf{Módulo testado} & \textbf{Cenários cobertos} \\
\midrule
\texttt{tb\_uart\_tx.v}      & \texttt{uart\_tx}      & Envio de 0x55, 0xAA, 0xFF, 0x00; tx\_ready \\
\texttt{tb\_lfsr8.v}         & \texttt{lfsr8}         & Nunca zero, período=255, freeze sem enable \\
\texttt{tb\_note\_track.v}   & \texttt{note\_track}   & Spawn, acerto (borda inf/centro/borda sup), escape, bad\_press, 3 slots, overflow, game\_active=0 \\
\texttt{tb\_packet\_sender.v}& \texttt{packet\_sender}& Todos os 20 bytes corretos, segundo pacote \\
\bottomrule
\end{tabularx}
\end{table}

\subsection{Executar com Icarus Verilog}

\begin{lstlisting}[style=shell]
cd /home/zocal/labdig_files/projeto_sem2

# Todos os testbenches
./run_tb.sh

# Um testbench especifico
./run_tb.sh tb_note_track
\end{lstlisting}

\subsection{Executar com ModelSim}

\begin{lstlisting}[style=shell]
# Compilar
vlog uart_tx.v lfsr8.v note_track.v tb_note_track.v

# Simular
vsim tb_note_track

# No console do ModelSim:
run -all
\end{lstlisting}

\subsection{Visualizar formas de onda}

Os testbenches geram arquivos \texttt{.vcd} que podem ser abertos no GTKWave:

\begin{lstlisting}[style=shell]
gtkwave tb_note_track.vcd &
\end{lstlisting}

Sinais recomendados para inspecionar em \texttt{tb\_note\_track.vcd}:
\begin{itemize}
    \item \texttt{frame\_tick} — pulsos de 60\,fps
    \item \texttt{note0}, \texttt{note1}, \texttt{note2} — ages das notas
    \item \texttt{hit\_pulse}, \texttt{bad\_press}, \texttt{escape\_pulse} — eventos de jogo
    \item \texttt{btn\_press} — pressão do botão
\end{itemize}

% ═══════════════════════════════════════════════════════════
\section{Resumo das Mudanças em Relação ao Semestre 1}
% ═══════════════════════════════════════════════════════════

\subsection{Problemas identificados no Semestre 1}

\begin{enumerate}
    \item \textbf{\texttt{fim\_tempo} e \texttt{perdeu} eram placeholders:}
    \begin{lstlisting}[style=verilog]
assign fim_tempo = botoes_raw[0] & botoes_raw[1]; // incorreto
assign perdeu    = botoes_raw[2] & botoes_raw[3]; // incorreto
    \end{lstlisting}
    Esses sinais não tinham lógica real — eram ativados acidentalmente ao pressionar dois botões simultâneos.

    \item \textbf{Contagem de tempo no \texttt{contador\_m}:}
    O parâmetro \texttt{M=3000} para 3 segundos implicava clock de 1\,kHz, provavelmente incorreto para a placa real.

    \item \textbf{Toda a lógica de jogo no Python:} notas, acertos, score, timer de 60\,s e condição de derrota estavam implementados em Python, contradizendo o objetivo do trabalho.
\end{enumerate}

\subsection{Soluções implementadas}

\begin{enumerate}
    \item \texttt{fim\_tempo} agora vem do \texttt{contador\_m} com $M=3600$ frames ($=60\,\text{s}$) no \texttt{fluxo\_dados}.
    \item \texttt{perdeu} agora é \texttt{miss\_reg >= 10}, calculado em hardware.
    \item \texttt{frame\_tick} gerado com $M=833\,333$ ciclos (50\,MHz/60\,fps).
    \item Quatro módulos novos implementam a lógica do jogo completamente em Verilog.
    \item Python reduzido a um renderer de $\approx 270$ linhas (era 470 linhas com toda a lógica).
\end{enumerate}

% ═══════════════════════════════════════════════════════════
\section*{Referências}
% ═══════════════════════════════════════════════════════════
\addcontentsline{toc}{section}{Referências}

\begin{enumerate}
    \item Intel/Altera. \textit{Cyclone V Device Handbook}. Disponível em: \url{https://www.intel.com/content/www/us/en/programmable/documentation/sam1403481054977.html}
    \item Terasic. \textit{DE1-SoC User Manual}. Disponível em: \url{https://www.terasic.com.tw/cgi-bin/page/archive.pl?Language=English&No=836&PartNo=4}
    \item Pygame. \textit{Pygame Documentation}. Disponível em: \url{https://www.pygame.org/docs/}
    \item pySerial. \textit{pySerial Documentation}. Disponível em: \url{https://pyserial.readthedocs.io/}
    \item Icarus Verilog. \textit{Icarus Verilog}. Disponível em: \url{http://iverilog.icarus.com/}
\end{enumerate}

\end{document}
